\section{Limits of scalar décalage and relative odd-disk envelopes}

\paragraph{Archive scope.}
The target-completion problem is the integral comparison problem of
cote 136, physical pp.~117, 119 and 121, archive footers 116, 118 and 120,
supplied scans 117, 119 and 121. The following explicit constructions
test possible repairs; they do not refute the original comparison statement.

\paragraph{ARCH-136-SCALAR-DECALAGE-NOGO.}
Let \(S=\Gamma_{\mathbb Z}(T)\), fix a positive weight \(m\), and consider
\[
 D=S\otimes\Lambda(a)\otimes\Gamma(b),\qquad
 |a|=1,\quad |b|=2,\quad da=b,\quad
 \operatorname{wt}(a)=\operatorname{wt}(b)=m .
\]
Choose nonzero integers \(s(q,w)\), and define the scalar décalage subcomplex
\[
 E^{w,q}=\{x\in s(q,w)D^{w,q}:
             dx\in s(q+1,w)D^{w,q+1}\}.
\]
Suppose \(E\) contains the full coefficient algebra \(S\) and is closed under its action and under the divided powers of \(D\). Then \(E\) cannot have vanishing positive cohomology in a weight tail in every fixed form degree.

\begin{proof}
The closed element \(s(2,m)b\) lies in \(E\). For every prime \(p\) and every \(r\geq0\), closure under divided powers and the coefficient action gives
\[
 s(2,m)^p\,T^{[r]}\gamma_p(b)\in E^{pm+r,2p}.
\]
The monomial \(T^{[r]}\gamma_p(b)\) is a primitive basis element of the corresponding free abelian group in \(D\). Consequently
\[
          s(2p,pm+r)\mid s(2,m)^p.                  \tag{136.DN1}
\]
Write \(w=pm+r\). In cochain degrees \(2p-1,2p\), the weight-\(w\) complex \(D\) is exactly
\[
 \mathbb Z\,T^{[r]}a\gamma_{p-1}(b)
       \xrightarrow{\ p\ }
 \mathbb Z\,T^{[r]}\gamma_p(b),
\]
with zero differential out of the right-hand term. Put
\(u=|s(2p-1,w)|\) and \(v=|s(2p,w)|\). Its scalar décalage has
\[
 H^{2p}(E^{w,*})
   =v\mathbb Z/
       \operatorname{lcm}(pu,v)\mathbb Z .
\]
Indeed the allowed domain elements are those multiples of \(u\) whose \(p\)-fold differential is divisible by \(v\). Vanishing of this quotient implies \(pu\mid v\), and in particular \(p\mid s(2p,w)\). If the required weight tail exists in degree \(2p\), take \(r\) large enough to enter it. Equation (136.DN1) then gives \(p\mid s(2,m)\). Since this holds for every prime, the nonzero integer \(s(2,m)\) is divisible by all primes, a contradiction.
\end{proof}

\paragraph{Exact scope.}
This excludes integer scalar divisibility décalages with arbitrary degree and weight dependence, already on one free odd disk. It does not exclude rational saturation, non-scalar kernels, quotient constructions, or algebra-dependent nonfree relations. It is distinct from the inherited AQ-136-FREE-ZERO-FORMS obstruction, which concerns simplicial full special-PD de Rham sources and a weight-one coefficient retraction.

\paragraph{ARCH-136-ODD-ENVELOPE-RELATIVE.}
An absolute positive cochain acyclic envelope of this disk does not, by that fact alone, give a relatively acyclic attachment. Here is an explicit integral example. Define an \(S\)-algebra \(A\) with even basis \(b_j=\gamma_j(b)\), \(j\geq0\), odd basis \(a_j\), \(j\geq1\), and
\[
 |b_j|=2j,\quad |a_j|=2j-1,\quad
 \operatorname{wt}(b_j)=\operatorname{wt}(a_j)=jm,
 \qquad da_j=b_j,\quad db_j=0.
\]
Set
\[
 b_i b_j=\binom{i+j}{i}b_{i+j},\qquad
 b_i a_j=\binom{i+j}{i}a_{i+j},\qquad
 a_i a_j=0.
\]
The even part has its usual divided powers. Then \(A\) is an integral strict differential PD algebra, free as an \(S\)-module, and \(S\to A\) is a cochain homotopy equivalence. The map \(D\to A\), \(a\mapsto a_1,b\mapsto b_1\), is strict PD. However the PD pushout along \(D\to S\), \(a,b\mapsto0\), has
\[
              H^3((A\otimes_D S)^{2m,*})
                        = \mathbb Z/2 .
\]
Here the displayed summand has coefficient weight zero; in higher total weights it is extended by the corresponding coefficient module.

\begin{proof}
The binomial identities establish associativity of the even action on the odd module. The differential is a derivation: on a product of two odd elements its right-hand side is
\(b_i a_j-a_i b_j=0\), and on an even--odd product the stated binomial formula gives the same result on both sides. Divided-power compatibility on the even part is immediate because all its elements are closed. The positive relative-weight complex is the direct sum of identity pairs \(Sa_j\to Sb_j\), so it contracts onto \(S\).

In the PD pushout, setting \(b_1=0\) sets all \(b_j=\gamma_j(b_1)\) to zero, and setting \(a_1=0\) gives the relation
\(2a_2=b_1a_1=0\). In relative weight two there is no other nonzero generator or further relation on \(a_2\), and its differential is zero. This gives the stated \(\mathbb Z/2\) class. The same point is visible in the free Moore repair
\[
 du=\gamma_2(b),\qquad dv=2u-ab,\qquad |u|=3,\ |v|=2:
\]
after \(a=b=0\), its relative weight-two complex is
\(\mathbb Z v\xrightarrow{2}\mathbb Z u\), with \(H^3=\mathbb Z/2\).
\end{proof}

\paragraph{Research gate.}
The relevant target problem is to choose all relative primitive relations coherently while preserving the fixed coefficient algebra and the degree-one period lattice. The examples above do not prove that such choices are impossible. In particular, the integer scalar no-go does not apply to a nonfree integral lattice inside a larger rational common target.

